\documentclass[11pt,a4paper]{article}
\usepackage[utf8]{inputenc}
\usepackage[T1]{fontenc}
\usepackage[french]{babel}
\frenchbsetup{StandardLists=true}
\usepackage{enumitem}
\usepackage{amssymb}
\usepackage{amsmath}
\usepackage[left=2cm,right=2cm,top=2cm,bottom=2cm]{geometry}
\usepackage[dvipsnames]{xcolor}
\usepackage{fouriernc}
\usepackage{setspace}
\singlespacing
\usepackage{fancyhdr}
\setlength{\headheight}{14pt}
\pagestyle{fancy}
\renewcommand\headrulewidth{1pt}
\fancyhead[L]{Lycée Denis Diderot}
\fancyhead[R]{Classe de 1BTS}
\setcounter{page}{1}\fancyfoot[C]{page \thepage}
\fancyfoot[R]{Année 2025-2026}
\fancyfoot[L]{Alexis RUIZ}
\usepackage{tikz}
\usepackage{multirow,tabularx,booktabs}
\usepackage{colortbl}
\usepackage[np]{numprint}
\usepackage{tcolorbox}
\usepackage{pifont}

\newcommand{\cadreblanc}[1]{%
\medskip\framebox[\linewidth][c]{%
\rule{0mm}{#1mm}}\par
\medskip}

\newcounter{numexos}
\setcounter{numexos}{0}
\newcommand{\bex}[2]{%
\def\bextitle{#1}%
\addtocounter{numexos}{1}%
\begin{tcolorbox}[colback=white,colframe=black!30,arc=3pt,boxrule=0.6pt,
  left=4mm,right=4mm,top=1mm,bottom=2mm]%
  \textbf{Exercice~\thenumexos\ifx\bextitle\empty\else~--~~#1\fi}\par
  #2%
\end{tcolorbox}}

\newcolumntype{L}[1]{>{\raggedright\arraybackslash}p{#1}}
\renewcommand{\arraystretch}{1.6}
\setlength{\parindent}{0pt}

\newtcolorbox{titrebox}{colback=white,colframe=black,arc=3pt,boxrule=1.5pt}

\begin{document}

\textbf{NOM :} \underline{\hspace{6cm}} \hfill \textbf{Classe :} \underline{\hspace{3cm}}\\[2mm]
\textbf{Prénom :} \underline{\hspace{5.5cm}}\\[4mm]
\hrule
\vspace{3mm}

\begin{titrebox}
\begin{center}
\textbf{\LARGE Contrôle -- Probabilités conditionnelles}\\[2mm]
\large Durée : 45 minutes \hfill Barème : \textbf{/10}
\end{center}
\end{titrebox}

\vspace{4mm}

%------------------------------------------------------------------
\bex{Tableau de probabilités \hfill \normalfont\textit{6 points}}{%
\smallskip
Une entreprise fabrique des écrans plats sur deux chaînes de production~:
\begin{itemize}
  \item $70\,\%$ des écrans sont produits par la chaîne $C_1$ et $30\,\%$ par la chaîne $C_2$.
  \item Parmi les écrans de $C_1$, $2\,\%$ présentent un défaut d'affichage.
  \item Parmi les écrans de $C_2$, $3{,}5\,\%$ présentent un défaut d'affichage.
\end{itemize}
On note $D$ l'événement \og l'écran présente un défaut d'affichage \fg. On prélève $\np{10000}$ écrans.
}

\textbf{1.} Compléter le tableau ci-dessous. \hfill \textit{(2 pts)}

\begin{center}
\renewcommand{\arraystretch}{2.2}
\begin{tabular}{|L{4.5cm}|c|c|c|}
\hline
 & Chaîne $C_1$ & Chaîne $C_2$ & Total \\
\hline
Avec défaut $D$ & \hspace{2.5cm} & \hspace{2.5cm} & \\
\hline
Sans défaut $\bar{D}$ & & & \\
\hline
Total & & & $\np{10000}$ \\
\hline
\end{tabular}
\end{center}

\textbf{2.} En déduire $P(D)$. \hfill \textit{(1 pt)}

\cadreblanc{15}

\textbf{3.} Définir l'événement $C_1 \cap D$ et calculer $P(C_1 \cap D)$. \hfill \textit{(1 pt)}

\cadreblanc{18}

\textbf{4.} Sachant qu'un écran est défectueux, calculer la probabilité qu'il provienne de la chaîne $C_1$.
On pourra utiliser la formule $P_D(C_1)=\dfrac{P(C_1 \cap D)}{P(D)}$. \hfill \textit{(2 pts)}

\cadreblanc{22}

\vspace{4mm}

%------------------------------------------------------------------
\bex{Indépendance de deux événements \hfill \normalfont\textit{4 points}}{%
\smallskip
On examine $\np{1000}$ pièces sorties d'une ligne de production pour deux types de défauts~: un défaut de \textbf{forme} ($F$) et un défaut de \textbf{couleur} ($C$).

\begin{center}
\renewcommand{\arraystretch}{1.8}
\begin{tabular}{|L{5cm}|c|c|c|}
\hline
 & Défaut $F$ & Pas de défaut $F$ & Total \\
\hline
Défaut $C$ & 10 & 70 & 80 \\
\hline
Pas de défaut $C$ & 40 & 880 & 920 \\
\hline
Total & 50 & 950 & $\np{1000}$ \\
\hline
\end{tabular}
\end{center}
}

\textbf{1.} Calculer $P(F)$ et $P(C)$. \hfill \textit{(1 pt)}

\cadreblanc{14}

\textbf{2.} Calculer $P(F \cap C)$. \hfill \textit{(1 pt)}

\cadreblanc{14}

\textbf{3.} Calculer $P(F) \times P(C)$. Comparer avec $P(F \cap C)$ et conclure quant à l'indépendance des événements $F$ et $C$. \hfill \textit{(2 pts)}

\cadreblanc{22}

\end{document}
